% Einheit 1 von 3 – Maßstab (bank/strahlensaetze/e1.jsonl, zusammenbau v0.1)
%% TODO Zweigzeile Teil 1: was hier gelernt wird (Satz in Schülersprache aus dem Einheitstitel) – Einheit 1
\einheitenkopf[e1]{Einheit 1 von 3 · Maßstab}
\zweigzeile{ab Kl. 5 · P10}
\setcounter{aufgabe}{11}

%% TODO Titel: Ich-kann-Satz fehlt in der Bank (Platzhalter: Kettenname) – Nr. 12
%% TODO Anweisung: ein Satz über der ersten Teilaufgabe fehlt in der Bank – Nr. 12
\begin{aufgabe}{Gleiche Einheit?}
\begin{teile}
\teil Plan: $4\,\text{cm}$, Wirklichkeit: $8\,\text{m}$ – gleiche Einheit? \kreuz{gleiche Einheit} \\ \kreuz{verschiedene Einheiten}
\end{teile}
\end{aufgabe}

%% TODO \verfahren{…}: Verfahrensüberschrift in Sachsprache fehlt in der Bank (2.3 b)
%% TODO Titel: Ich-kann-Satz fehlt in der Bank (Platzhalter: Kettenname) – Nr. 13
%% TODO Anweisung: ein Satz über der ersten Teilaufgabe fehlt in der Bank – Nr. 13
\begin{aufgabe}{Maßstab umrechnen}
%% TODO \rechenplatz{n}: Schrittzahl des Grundfalls fehlt in der Bank (nur bei fester Schreibform, 2.3 b)
\begin{teile}
\teil Maßstab $1 : 20$ – ist die Zeichnung kleiner oder größer als die Wirklichkeit? \kreuz{kleiner} \\ \kreuz{größer}
\teil Maßstab $1 : 20$ – wie viel ist $1\,\text{cm}$ auf dem Plan in Wirklichkeit? \leerfeld[cm]
\teil Ein Tisch ist $150\,\text{cm}$ lang. Maßstab $1 : 25$ – wie lang ist er auf dem Plan? \leerfeld[cm]
\teil Auf dem Plan ist ein Zimmer $5\,\text{cm}$ lang. Maßstab $1 : 80$ – wie lang ist es in Wirklichkeit? \leerfeld[m]
\teil Ein Garten ist $18\,\text{m}$ lang. Maßstab $1 : 300$ – wie lang ist er auf dem Plan? \leerfeld[cm]
\teil Ein Schrank ist $2{,}1\,\text{m}$ hoch und $1{,}5\,\text{m}$ breit. Maßstab $1 : 30$ – Höhe und Breite auf dem Plan? Höhe \leerfeld[cm] , Breite \leerfeld[cm]
\teil Ein Käfer ist $8\,\text{mm}$ lang. Zeichnung im Maßstab $4 : 1$ – wie lang ist er auf der Zeichnung? \leerfeld[mm]
\teil Auf dem Plan $5\,\text{cm}$, in Wirklichkeit $15\,\text{m}$ – welcher Maßstab? 1 : \leerfeld
\teil Karte im Maßstab $1 : 50\,000$, auf der Karte $6\,\text{cm}$ – wie viele Kilometer? \leerfeld[km]
\teil Ein Leuchtturm ist $28{,}0\,\text{m}$ hoch und hat unten einen Durchmesser von $5{,}6\,\text{m}$. Für eine Ausstellung wird ein Modell im Maßstab $1 : 40$ gebaut. Berechne die Höhe und den Durchmesser des Modells in Zentimetern. \hfill (P10 2018 OS) \\ Höhe \leerfeld[cm] , Durchmesser \leerfeld[cm]
\end{teile}
\end{aufgabe}

%% TODO \verfahren{…}: Verfahrensüberschrift in Sachsprache fehlt in der Bank (2.3 b)
%% TODO Titel: Ich-kann-Satz fehlt in der Bank (Platzhalter: Kettenname) – Nr. 14
%% TODO Anweisung: ein Satz über der ersten Teilaufgabe fehlt in der Bank – Nr. 14
\begin{aufgabe}{Maßstabsgerecht zeichnen}
%% TODO \rechenplatz{n}: Schrittzahl des Grundfalls fehlt in der Bank (nur bei fester Schreibform, 2.3 b)
\begin{teile}
\teil Ein Beet ist $4\,\text{m}$ lang und $2\,\text{m}$ breit. Das Zeichenfeld ist $15\,\text{cm}$ breit und $9\,\text{cm}$ hoch. Welcher Maßstab passt? \kreuz{$1 : 20$} \\ \kreuz{$1 : 40$} \\ \kreuz{$1 : 1\,000$}
\teil Im Raster steht ein Kästchen für $2\,\text{m}$. Wie viele Kästchen lang ist eine $14\,\text{m}$ lange Mauer? Trage die Strecke im Raster ab.

\begin{ksys}[xmin=0,xmax=10,ymin=0,ymax=3]\end{ksys}
\leerfeld[Kästchen]
\teil Das Rechteck $ABCD$ ist $4$ Kästchen breit und $3$ Kästchen hoch. Zeichne es doppelt so groß.

\begin{ksys}[xmin=0,xmax=14,ymin=0,ymax=10]\punkt{1}{1}{A}\punkt{5}{1}{B}\punkt{5}{4}{C}\punkt{1}{4}{D}\end{ksys}
\teil Ein Spielfeld ist $18\,\text{m}$ lang und $9\,\text{m}$ breit. Zeichne es im Maßstab $1 : 300$ und beschrifte es mit den wirklichen Maßen.

\begin{ksys}[xmin=0,xmax=10,ymin=0,ymax=6]\end{ksys}
\teil Ein runder Teich hat einen Durchmesser von $6\,\text{m}$. Zeichne ihn im Maßstab $1 : 100$ mit dem Zirkel.

\begin{ksys}[xmin=0,xmax=10,ymin=0,ymax=6]\end{ksys}
\teil Eine zylinderförmige Säule steht mittig auf einem würfelförmigen Sockel. Welche Figur ist die Draufsicht? \kreuz{Kreis im Quadrat} \\ \kreuz{Rechteck auf Rechteck} \\ \kreuz{Kreis neben Quadrat}
\teil Ein runder Tisch mit $1{,}6\,\text{m}$ Durchmesser steht mitten auf einem quadratischen Teppich mit $2{,}4\,\text{m}$ Seitenlänge. Zeichne die Draufsicht im Maßstab $1 : 40$.

\begin{ksys}[xmin=0,xmax=10,ymin=0,ymax=6]\end{ksys}
\teil Ein Carport ist $6\,\text{m}$ lang und $3\,\text{m}$ breit. Das Zeichenfeld ist $15\,\text{cm}$ breit und $9\,\text{cm}$ hoch. Wähle einen passenden Maßstab, gib ihn an, zeichne und beschrifte.

\begin{ksys}[xmin=0,xmax=10,ymin=0,ymax=6]\end{ksys}
\teil In einer Zeichnung im Maßstab $1 : 20$ ist ein Schrank $4{,}5\,\text{cm}$ breit, die Nische für ihn $4{,}2\,\text{cm}$. Passt der Schrank in die Nische? Wie viel fehlt oder bleibt in Wirklichkeit?
\teil Eine zylinderförmige Futtertonne ist $90\,\text{cm}$ hoch und hat den Radius $33\,\text{cm}$. Sie soll mit einer rechteckigen Holzplatte abgedeckt werden, die $60\,\text{cm}$ breit und $85\,\text{cm}$ lang ist. Zeichne die Draufsicht von Tonne und Platte in einem passenden Maßstab in das Karofeld, gib den Maßstab an und beschrifte. Entscheide, ob die Platte die Tonne vollständig abdeckt. \hfill (P10 2021 OS) \\

\begin{ksys}[xmin=0,xmax=13,ymin=0,ymax=14]\end{ksys}
Maßstab: \leerfeld
\end{teile}
\end{aufgabe}

%% TODO Titel: Ich-kann-Satz fehlt in der Bank (Platzhalter: Kettenname) – Nr. 15
%% TODO Anweisung: ein Satz über der ersten Teilaufgabe fehlt in der Bank – Nr. 15
\begin{aufgabe}{Maßstab umrechnen – Fehler finden, Begründen, Anwendung, Darstellungswechsel}
\begin{teile}
\teil Mia soll eine $9\,\text{m}$ lange Hecke im Maßstab $1 : 30$ zeichnen. Sie rechnet: \rechnung{900\,\text{cm} \cdot 30 &= 27\,000\,\text{cm}} Finde den Fehler und rechne richtig.
\teil Warum wird eine Zeichnung im Maßstab $1 : 80$ kleiner als eine im Maßstab $1 : 40$? Begründe.
\teil Auf einer Wanderkarte im Maßstab $1 : 50\,000$ ist die Route $16\,\text{cm}$ lang. Die Familie schafft $4\,\text{km}$ in der Stunde. Wie lange ist sie unterwegs? \leerfeld Stunden
\teil Die Maßstabsleiste einer Karte zeigt: $2\,\text{cm}$ ≙ $1\,\text{km}$. Schreibe den Maßstab als $1 : n$. 1 : \leerfeld
\end{teile}
\end{aufgabe}
